\documentclass[11pt]{article}
\usepackage[utf8]{inputenc}
\usepackage{amsmath}
\usepackage{amsfonts}
\usepackage{amssymb}
\usepackage{bm}
\usepackage{geometry}
\usepackage{xcolor}
\usepackage{booktabs}
\usepackage{multirow}
\usepackage{algorithm}
\usepackage{algorithmic}

\geometry{margin=0.8in}

\title{Synthetic Non-Stationary Spiking Data Generation:\\
Dale-Law-Constrained Recurrent Networks with\\
a Switching Poisson GLM}
\author{Jan Balewski, NERSC}
\date{\today}

\begin{document}

\maketitle

\begin{abstract}
We describe the two-script pipeline that produces synthetic non-stationary
neural population spike trains for benchmarking latent-state inference
algorithms (e.g.\ PRISM-EM).
\textbf{Stage~1} (\texttt{gen\_daleMatrices3.py}) constructs a single
Dale-law-compliant weight matrix $A_{\text{true}}$ with a prescribed
spectral radius $R<1$ and $M$ bias vectors $\{B_m\}_{m=0}^{M-1}$,
one per dynamical state.
Stationary spike trains under each bias are generated as a
firing-rate sanity check.
\textbf{Stage~2} (\texttt{gen\_nonStationarySpikes3.py}) loads the
ground-truth dictionary, schedules transitions among the $M$ states,
tracks a smoothly evolving simplex coefficient $c_t$, and generates
the non-stationary spike train from a Poisson GLM with a
time-varying effective bias $B_{\mathrm{eff},t} = \sum_m c_{t,m} B_m$.
An oracle state decoder (exact Poisson log-likelihood) is also
computed and saved as an upper bound for any inference algorithm.
\end{abstract}

\tableofcontents
\newpage

% ===============================================================
\section{Overview}
\label{sec:overview}
% ===============================================================

The pipeline produces ground-truth data for a network of $N$ neurons
divided into $N_E$ excitatory and $N_I = N - N_E$ inhibitory cells.
The network obeys Dale's principle: each neuron is either purely
excitatory or purely inhibitory.
All neurons share one connectivity matrix $A_{\text{true}}$; what
differs across $M$ states is only the bias (log-baseline firing-rate)
vector $B_m$.

\paragraph{Stage 1 (Section~\ref{sec:dale_matrices}).}
\texttt{gen\_daleMatrices3.py} builds and saves:
\begin{itemize}
  \item a sparse binary connectivity mask $E_{\text{true}} \in
        \{0,1\}^{N \times N}$;
  \item a single weight matrix $A_{\text{true}}$ with
        $\rho(A_{\text{true}}) = R$;
  \item $M$ bias vectors $B_{\text{true}} \in \mathbb{R}^{M \times N}$,
        one per state.
\end{itemize}
Stationary spikes under each $B_m$ are simulated and their
firing-rate statistics are reported.

\paragraph{Stage 2 (Section~\ref{sec:nonstat_spikes}).}
\texttt{gen\_nonStationarySpikes3.py} loads the Stage-1 output and:
\begin{itemize}
  \item draws a target-state sequence $S_{\text{true}}[t]$
        via a Markov chain or round-robin schedule;
  \item maintains a simplex coefficient vector $c_t \in \Delta^{M-1}$
        that blends the $M$ bias vectors smoothly;
  \item samples $Y_t \sim \mathrm{Poisson}$ from the resulting
        time-varying Poisson GLM;
  \item computes the oracle state $S_{\text{oracle}}[t]$ by
        maximum Poisson log-likelihood.
\end{itemize}

\paragraph{Notation.}
$N$: total neurons; $N_E$: excitatory count (default $\lfloor 0.8N \rfloor$);
$N_I = N - N_E$: inhibitory count.
$M$: number of states (= \texttt{len(Boffsets)}).
$\Delta t$: time-bin width in seconds (default $0.01$\,s).
$T$: number of time bins.
$R$: target spectral radius (default $0.3$).
$\eta_{\mathrm{clip}} = 5$: overflow-prevention threshold
($\approx [10^{-3}, 10^3]$\,Hz effective rate range).
Column convention: $(A_{\text{true}})_{ij}$ is the synaptic weight
\emph{from} presynaptic neuron $j$ \emph{to} postsynaptic neuron $i$.

% ===============================================================
\section{Stage 1: Dale Matrix Generation}
\label{sec:dale_matrices}
% ===============================================================

\subsection{Command-Line Parameters}

\begin{table}[h]
\centering
\begin{tabular}{lll}
\toprule
CLI flag & Default & Description \\
\midrule
\texttt{--num\_neurons} $N$ & 50 & Total neuron count \\
\texttt{--num\_excite} $N_E$ & $\lfloor 0.8N \rfloor$ &
  Excitatory neuron count \\
\texttt{--edge\_prob} $[p_\ell, p_h]$ & $[0.05,\, 0.20]$ &
  Per-neuron connectivity probability range \\
\texttt{--spectral\_radius} $R$ & 0.3 &
  Target spectral radius \\
\texttt{--idleRate} $[f_{\min}, f_{\max}]$ & $[15,\, 30]$\,Hz &
  Idle firing-rate band \\
\texttt{--Boffsets} $\{\delta_m\}$ & $[0.0]$ &
  Per-state log-rate offsets (one bias vector per entry) \\
\texttt{--num\_steps} $T$ & 10\,001 &
  Steps for the stationary sanity-check simulation \\
\texttt{--step\_size} $\Delta t$ & 0.01\,s & Time-bin width \\
\texttt{--basePath} & \emph{(scratch dir)} &
  Root output directory \\
\bottomrule
\end{tabular}
\caption{Key command-line arguments of \texttt{gen\_daleMatrices3.py}.}
\label{tab:dale_args}
\end{table}

\subsection{Algorithm}

\begin{algorithm}[H]
\caption{Dale Matrix Generation (\texttt{gen\_daleMatrices3.py})}
\begin{algorithmic}[1]
\STATE \textbf{Step 1:} Generate sparse binary connectivity mask
       $E_{\text{true}} \in \{0,1\}^{N\times N}$
       (\texttt{generate\_sparse\_mask})
\STATE \textbf{Step 2:} Generate weight matrix $A_{\text{true}}$ with
       Dale E/I balance, negative self-loops, and spectral radius $R$
       (\texttt{init\_W})
\STATE \textbf{Step 3:} For each offset $\delta_m$ in \texttt{Boffsets}:
\begin{ALC@g}
  \STATE Generate bias vector $B_m$ (\texttt{set\_flat\_selfSpiking})
  \STATE Simulate stationary spikes $Y^{(m)}$
         (\texttt{gen\_stationary\_lag1\_poisson})
  \STATE Compute firing-rate statistics (\texttt{estimate\_rates})
\end{ALC@g}
\STATE \textbf{Step 4:} Save ground-truth and spike files
\end{algorithmic}
\end{algorithm}

\subsubsection{Step 1: Sparse Connectivity Mask}
\label{sec:mask}

Each neuron $i$ independently draws its connection probability:
\begin{equation}
  q_i \;\sim\; \mathrm{Uniform}(p_\ell,\; p_h).
\end{equation}
The binary mask entry is
\begin{equation}
  (E_{\text{true}})_{ij} =
  \begin{cases}
    1 & \text{if } u_{ij} < q_i, \quad
        u_{ij} \sim \mathrm{Uniform}(0,1) \text{ i.i.d.,} \\
    0 & \text{otherwise.}
  \end{cases}
\end{equation}
Self-loops are \emph{always} included: $(E_{\text{true}})_{ii} = 1$
for all $i$, encoding the mandatory negative self-connection added below.

\subsubsection{Step 2: Weight Matrix with Dale's Principle}
\label{sec:weights}

\paragraph{E/I balance ratio.}
Let $r = N_E / N_I$.
This ratio is used to set the mean inhibitory weight magnitude so
that the expected signed row-sum of the full matrix is
approximately zero.

\paragraph{Raw weight assignment (row-based, weight spread $v = 0.2$).}
\begin{equation}
  W_{ij} =
  \begin{cases}
    \mathrm{Uniform}(1-v,\; 1+v) &
      i \in \{0,\dots,N_E-1\} \quad (\text{excitatory rows}), \\[4pt]
    \mathrm{Uniform}(-r(1+v),\; -r(1-v)) &
      i \in \{N_E,\dots,N-1\} \quad (\text{inhibitory rows}).
  \end{cases}
  \label{eq:raw_weights}
\end{equation}
Because the \emph{row} index $i$ determines the neuron's type, every
outgoing weight of an excitatory neuron is positive and every
outgoing weight of an inhibitory neuron is negative.  This is
Dale's principle.

\paragraph{Applying the connectivity mask.}
\begin{equation}
  W \;\leftarrow\; W \odot E_{\text{true}}.
\end{equation}

\paragraph{Enforcing negative self-loops.}
Any positive diagonal entry (excitatory self-loop) is sign-flipped:
\begin{equation}
  W_{ii} \;\leftarrow\; -|W_{ii}| \quad \forall i,
\end{equation}
so that all self-connections are inhibitory (stabilising).

\paragraph{Spectral normalization.}
Let $\rho_0 = \max_k |\lambda_k(W)|$ be the current spectral radius.
The connectivity matrix is set to
\begin{equation}
  A_{\text{true}} \;\leftarrow\; \frac{R}{\rho_0}\, W,
  \label{eq:spectral}
\end{equation}
so $\rho(A_{\text{true}}) = R$ exactly.
Rescaling preserves all signs, the zero pattern, Dale's law, and
the negative self-loops.
For $R < 1$ the operator $A_{\text{true}}$ is contracting, which is
a sufficient condition for the stationary Poisson VAR(1) to have a
well-defined stationary distribution.

\subsubsection{Step 3: Bias Vectors and Stationary Spike Check}
\label{sec:bias}

\paragraph{Bias generation (\texttt{set\_flat\_selfSpiking}).}
A size-scaling factor $s = \sqrt{N/100}$ adjusts the excitatory
correction for network size.
For each state offset $\delta_m$ (index $m = 0,\dots,M-1$),
independent draws
\begin{equation}
  B_m^{(i)} \;\sim\;
  \mathrm{Uniform}\!\bigl(\ln(R\,f_{\min}),\;\ln(R\,f_{\max})\bigr),
  \quad i = 0,\dots,N-1,
\end{equation}
followed by state- and type-dependent shifts:
\begin{equation}
  B_m^{(i)} \;\leftarrow\;
  \begin{cases}
    B_m^{(i)} + 1 + \delta_m - 1.5R - s
      & i < N_E \quad (\text{excitatory}), \\[4pt]
    B_m^{(i)} + \delta_m
      & i \geq N_E \quad (\text{inhibitory}).
  \end{cases}
  \label{eq:bias}
\end{equation}
The excitatory correction $-(1.5R + s - 1)$ lowers the baseline
excitatory rate relative to inhibitory neurons, preventing runaway
excitation.
Larger $\delta_m$ uniformly raises activity across the network,
creating distinct mean-rate states.

\paragraph{Stationary spike simulation (evaluation only).}
For each bias vector $B_m$ a lag-1 Poisson VAR(1) process is
simulated over $T$ bins:
\begin{equation}
  \eta_t = A_{\text{true}}\,Y_{t-1} + B_m,
  \qquad
  \tilde{\eta}_t = \min(\eta_t,\;\eta_{\mathrm{clip}}),
  \qquad
  \lambda_t = \exp(\tilde{\eta}_t)\,\Delta t,
  \qquad
  Y_t \;\sim\; \mathrm{Poisson}(\lambda_t),
  \label{eq:stationary}
\end{equation}
with initial condition $Y_0 \sim \mathrm{Poisson}(\exp(B_m)\,\Delta t)$.
These stationary spike arrays (shape $(T, N)$, one per bias vector)
are used \emph{only} to verify that the firing rates are in a
plausible range; they are \emph{not} the final non-stationary output.

Firing-rate statistics computed via \texttt{estimate\_rates}
(\texttt{UtilDalePoisson.py}) include:
\begin{itemize}
  \item mean per-neuron firing rates (Hz),
  \item rate variance over non-overlapping $T_{\mathrm{var}} = 5$\,s windows (Hz$^2$),
  \item Fano factors ($\mathrm{Var}[\text{count}] / \mathrm{Mean}[\text{count}]$),
  \item consecutive cross-neuron coincidence rate (Hz per neuron).
\end{itemize}

\subsection{Output Files (Stage 1)}

All files are written to \texttt{<basePath>/truthDale/}.

\begin{table}[h]
\centering
\begin{tabular}{llll}
\toprule
File & Array key & Shape & Description \\
\midrule
\multirow{3}{*}{\texttt{<dataName>.simTruth.npz}}
  & \texttt{A\_true} & $(N,N)$ float &
      Shared weight matrix, $\rho(A_{\text{true}})=R$ \\
  & \texttt{B\_true} & $(M,N)$ float &
      Per-state bias vectors \\
  & \texttt{E\_true} & $(N,N)$ int &
      Binary connectivity mask \\
\midrule
\multirow{4}{*}{\texttt{<dataName>.spikes.npz}}
  & \texttt{spikes} & $(M,T,N)$ uint8 &
      Stationary spikes per state (sanity check) \\
  & \texttt{single\_rates} & $(M,N)$ float &
      Mean firing rates (Hz) per state \\
  & \texttt{sigle\_rates\_var} & $(M,N)$ float &
      Rate variance (Hz$^2$) per state \\
  & \texttt{single\_fano\_fact} & $(M,N)$ float &
      Fano factors per state \\
\bottomrule
\end{tabular}
\caption{Output arrays of Stage~1.
  $M = \texttt{len(Boffsets)}$; $T = \texttt{num\_steps}$.}
\label{tab:stage1_output}
\end{table}

The metadata dictionary \texttt{dale\_conf} records $N$, $N_E$, $R$,
edge-probability range, idle-rate range, and offset list.
The metadata dictionary \texttt{evol\_conf} records $T$, $\Delta t$,
total simulated time, and $\eta_{\mathrm{clip}}$.

% ===============================================================
\section{Stage 2: Non-Stationary Spike Generation}
\label{sec:nonstat_spikes}
% ===============================================================

\subsection{Command-Line Parameters}

\begin{table}[h]
\centering
\begin{tabular}{lll}
\toprule
CLI flag & Default & Description \\
\midrule
\texttt{--inputStates} & --- &
  Base name of the \texttt{.simTruth.npz} file (required) \\
\texttt{--num\_steps} $T$ & from \texttt{evol\_conf} &
  Number of time bins \\
\texttt{--true\_dwell\_sec} $t_d$ & 1.0\,s &
  Mean dwell time per state \\
\texttt{--state\_change\_speed} $\nu$ & 0.33 &
  Max simplex-coefficient change per bin \\
\texttt{--schedule} & \texttt{mc} &
  State schedule: Markov chain (\texttt{mc}) or round-robin (\texttt{rr}) \\
\texttt{--seed} & 42 & Random seed \\
\bottomrule
\end{tabular}
\caption{Key command-line arguments of \texttt{gen\_nonStationarySpikes3.py}.}
\label{tab:nonstat_args}
\end{table}

\subsection{Algorithm}

\begin{algorithm}[H]
\caption{Non-Stationary Spike Generation
  (\texttt{gen\_nonStationarySpikes3.py})}
\begin{algorithmic}[1]
\STATE Load $(A_{\text{true}}, B_{\text{true}})$ from
       \texttt{<inputStates>.simTruth.npz}
\STATE Build target-state sequence $S_{\text{true}}[t]$
       via selected schedule (Section~\ref{sec:scheduling})
\STATE Simulate switching spikes and record simplex coefficients $C_{\text{true}}$
       (\texttt{simulate\_switching\_poisson}, Section~\ref{sec:switching})
\STATE Compute oracle state sequence $S_{\text{oracle}}$
       (\texttt{compute\_oracle\_states\_comA}, Section~\ref{sec:oracle})
\STATE Compute firing-rate statistics (\texttt{estimate\_rates})
\STATE Save spike and ground-truth files
\end{algorithmic}
\end{algorithm}

\subsection{State Scheduling}
\label{sec:scheduling}

The mean dwell time in bins is
$\tau_d = \lceil t_d / \Delta t \rceil$.

\paragraph{Markov chain (\texttt{mc}, default).}
A discrete-time Markov chain with geometric dwell times.
The exit probability per bin is $p_{\mathrm{exit}} = 1/\tau_d$;
on exit the chain transitions uniformly to one of the $M-1$ other states.
The symmetric $M \times M$ transition matrix stored in the output is:
\begin{equation}
  T_{sm} =
  \begin{cases}
    1 - p_{\mathrm{exit}} & s = m \quad (\text{self-transition}), \\
    p_{\mathrm{exit}}/(M-1) & s \neq m \quad (\text{cross-transition}).
  \end{cases}
  \label{eq:trans}
\end{equation}
A minimum dwell of
$\tau_{\min} = \lceil 0.3\,\tau_d \rceil$ bins is enforced
to prevent spuriously short visits.
State coverage across the full record is uncontrolled and subject
to random fluctuations.

\paragraph{Round-robin (\texttt{rr}).}
States cycle deterministically as $0, 1, \dots, M-1, 0, 1, \dots$
Each visit lasts exactly $\tau_d$ bins (the last visit in the
record is trimmed to fit $T$).
This guarantees that all states receive equal coverage
($T/M \pm \tau_d$ bins each), eliminating sampling variability.

\subsection{Smooth Switching Poisson Simulation}
\label{sec:switching}

\paragraph{Simplex coefficient vector.}
A vector $c_t \in \Delta^{M-1}$
(non-negative entries summing to 1) tracks the current mixture of
states.
The one-hot target for the scheduled state at bin $t$ is
$e_m = \mathbf{1}[m = S_{\text{true}}[t]]$.
At each bin, $c_t$ is updated by a clipped gradient step and
renormalized:
\begin{align}
  c_t &\;\leftarrow\; c_{t-1}
        + \mathrm{clip}\!\bigl(e - c_{t-1},\; -\nu,\; \nu\bigr), \\
  c_t &\;\leftarrow\; c_t \;/\; \|c_t\|_1.
  \label{eq:smooth_update}
\end{align}
The speed parameter $\nu \in (0,1]$ controls how quickly the mixture
tracks the target: $\nu = 1$ gives instantaneous switching;
small $\nu$ produces gradual multi-bin transitions.
The transition duration is approximately $1/\nu$ bins.

\paragraph{Effective bias.}
The convex mixture of per-state bias vectors is
\begin{equation}
  B_{\mathrm{eff},t}
  = \sum_{m=0}^{M-1} c_{t,m}\, B_m
  = \sum_{m=0}^{M-1} c_{t,m}\, B_{\text{true}}[m,:].
  \label{eq:Beff}
\end{equation}
This interpolates smoothly between the per-state firing-rate
baselines as the state label switches.

\paragraph{Non-stationary spike generation.}
At every bin $t = 0, \dots, T-1$:
\begin{equation}
  \eta_t = A_{\text{true}}\,Y_{t-1} + B_{\mathrm{eff},t},
  \qquad
  \tilde{\eta}_t = \min(\eta_t,\;\eta_{\mathrm{clip}}),
  \qquad
  \lambda_t = \exp(\tilde{\eta}_t)\,\Delta t,
  \qquad
  Y_t \;\sim\; \mathrm{Poisson}(\lambda_t).
  \label{eq:nonstat}
\end{equation}
This is the same Poisson VAR(1) model as in Eq.~\eqref{eq:stationary},
except that the bias $B_{\mathrm{eff},t}$ is now time-varying.
The connectivity matrix $A_{\text{true}}$ is unchanged across all states.

\subsection{Oracle State Decoder}
\label{sec:oracle}

An oracle decoder assigns the most likely state to each bin using
the \emph{exact} ground-truth parameters $(A_{\text{true}}, \{B_m\})$:
\begin{equation}
  S_{\text{oracle}}[t]
  = \arg\max_{m \in \{0,\dots,M-1\}}\;
    \sum_{j=0}^{N-1}
    \Bigl[
      Y_{t,j}\,\tilde{\eta}_{t,j}^{(m)}
      - \exp\!\bigl(\tilde{\eta}_{t,j}^{(m)}\bigr)\,\Delta t
    \Bigr],
  \label{eq:oracle}
\end{equation}
where
$\eta_{t,j}^{(m)} = \bigl(A_{\text{true}}\,Y_{t-1}\bigr)_j + B_{m,j}$
and $\tilde{\eta}^{(m)} = \min(\eta^{(m)}, \eta_{\mathrm{clip}})$.
This is the exact Poisson log-likelihood score for state $m$ given
the observed spike vector $Y_t$, computed independently at each bin.
The oracle accuracy is an upper bound for any latent-state inference
algorithm.

Per-state oracle scores and overall accuracy are printed at the end of
the run and stored in the \texttt{oracle\_eval} sub-dictionary of the
metadata.

\subsection{Output Files (Stage 2)}

All files are written to \texttt{<basePath>/spikesData/}.

\begin{table}[h]
\centering
\begin{tabular}{llll}
\toprule
File & Array key & Shape & Description \\
\midrule
\multirow{2}{*}{\texttt{<dataName>.spikes.npz}}
  & \texttt{spikes} & $(T,N)$ int32 &
      Non-stationary spike counts \\
  & \texttt{single\_rates} & $(N,)$ float &
      Mean firing rates over the full record (Hz) \\
\midrule
\multirow{7}{*}{\texttt{<dataName>.prismTruth.npz}}
  & \texttt{A\_true} & $(N,N)$ float &
      Shared weight matrix (copy from Stage 1) \\
  & \texttt{B\_true} & $(M,N)$ float &
      Per-state bias vectors (copy from Stage 1) \\
  & \texttt{S\_true} & $(T,)$ int32 &
      Target state index at each bin \\
  & \texttt{C\_true} & $(T,M)$ float32 &
      Simplex coefficient vector at each bin \\
  & \texttt{S\_oracle} & $(T,)$ int32 &
      Oracle (max-likelihood) state at each bin \\
  & \texttt{state\_transition} & $(M,M)$ float32 &
      Transition matrix used/implied by schedule \\
  & \texttt{sigle\_rates\_var} & $(N,)$ float &
      Per-neuron rate variance (Hz$^2$) \\
  & \texttt{single\_fano\_fact} & $(N,)$ float &
      Per-neuron Fano factors \\
\bottomrule
\end{tabular}
\caption{Output arrays of Stage~2.
  $T = \texttt{num\_steps}$; $M = \texttt{num\_states}$.}
\label{tab:stage2_output}
\end{table}

The metadata dictionary \texttt{evol\_conf} records $T$, $\Delta t$,
total time, $\nu$, requested and effective dwell times (seconds and bins),
$M$, random seed, and schedule type.
The \texttt{oracle\_eval} sub-dictionary records the overall oracle
accuracy and per-state breakdown.

% ===============================================================
\section{Notation Summary}
\label{sec:notation}
% ===============================================================

\begin{table}[h]
\centering
\renewcommand{\arraystretch}{1.2}
\begin{tabular}{lll}
\toprule
Symbol & Domain & Description \\
\midrule
$N$ & $\mathbb{Z}_{>0}$ & Total number of neurons \\
$N_E$ & $\mathbb{Z}_{>0}$ & Excitatory neuron count (rows $0\dots N_E-1$) \\
$N_I = N-N_E$ & $\mathbb{Z}_{>0}$ & Inhibitory neuron count \\
$M$ & $\mathbb{Z}_{>0}$ &
  Number of states ($= \texttt{len(Boffsets)}$) \\
$T$ & $\mathbb{Z}_{>0}$ & Number of time bins \\
$\Delta t$ & $\mathbb{R}_{>0}$ & Time-bin width (s); default $0.01$ \\
$R$ & $(0,1)$ & Target spectral radius \\
$r = N_E/N_I$ & $\mathbb{R}_{>0}$ & E/I balance ratio \\
$v = 0.2$ & --- & Weight spread (hard-coded) \\
$E_{\text{true}}$ & $\{0,1\}^{N\times N}$ & Binary connectivity mask \\
$A_{\text{true}}$ & $\mathbb{R}^{N\times N}$ &
  Shared weight matrix, $\rho(A_{\text{true}}) = R$ \\
$B_m$ & $\mathbb{R}^N$ & Bias (log-baseline rate) vector for state $m$ \\
$B_{\text{true}}$ & $\mathbb{R}^{M\times N}$ & Stacked per-state biases \\
$Y_t$ & $\mathbb{Z}_{\geq 0}^N$ & Spike-count vector at bin $t$ \\
$\eta_{\mathrm{clip}}$ & $\mathbb{R}_{>0}$ &
  Rate clipping threshold; fixed at $5$ \\
\midrule
$S_{\text{true}}[t]$ & $\{0,\dots,M-1\}$ &
  Scheduled target state at bin $t$ \\
$c_t$ & $\Delta^{M-1}$ &
  Simplex coefficient vector at bin $t$ \\
$\nu$ & $(0,1]$ &
  Max coefficient change per bin (\texttt{state\_change\_speed}) \\
$B_{\mathrm{eff},t}$ & $\mathbb{R}^N$ &
  Effective bias: convex mixture $\sum_m c_{t,m} B_m$ \\
$\tau_d$ & $\mathbb{Z}_{>0}$ &
  Mean dwell time in bins \\
$S_{\text{oracle}}[t]$ & $\{0,\dots,M-1\}$ &
  Oracle max-likelihood state at bin $t$ \\
\bottomrule
\end{tabular}
\caption{Notation used throughout this document.}
\label{tab:notation}
\end{table}

% ===============================================================
\section{Full Pipeline and Suggested Commands}
\label{sec:pipeline}
% ===============================================================

\begin{enumerate}

\item \textbf{Generate ground-truth matrices} (Stage 1):
\begin{verbatim}
./gen_daleMatrices3.py --num_neurons 100 --num_excite 70 \
    --spectral_radius 0.5 --Boffsets 0 10 20 --dataName myData \
    --basePath $basePath
\end{verbatim}
Produces \texttt{myData.simTruth.npz} ($A_{\text{true}}$,
$B_{\text{true}}$, $E_{\text{true}}$; $M=3$ states here) and
\texttt{myData.spikes.npz} (stationary sanity-check arrays,
shape $(3, T, N)$).

\item \textbf{Inspect matrices and stationary spikes}:
\begin{verbatim}
./view_daleMatrix3.py  --basePath $basePath --dataName myData \
    -p b -m 0 -X -p a c d
./view_spikesTrain3.py --basePath $basePath --dataName myData \
    --time_range_sec 0 20 -p b -m 0 -X
\end{verbatim}

\item \textbf{Generate non-stationary spike train} (Stage 2):
\begin{verbatim}
./gen_nonStationarySpikes3.py --basePath $basePath \
    --inputStates myData --true_dwell_sec 1.0
\end{verbatim}
Produces \texttt{<dataName>.spikes.npz} (shape $(T, N)$) and
\texttt{<dataName>.prismTruth.npz} ($S_{\text{true}}$,
$C_{\text{true}}$, $S_{\text{oracle}}$, transition matrix, and
copies of $A_{\text{true}}$, $B_{\text{true}}$).

\item \textbf{Inspect non-stationary spikes}:
\begin{verbatim}
./view_spikesTrain3.py --basePath $basePath \
    --dataName <dataName> --idxState -1 \
    --time_range_sec 0 15 -p b
\end{verbatim}

\item \textbf{Fit with PRISM-EM}:
\begin{verbatim}
./fitPrismEM.sh --basePath $basePath --dataName <dataName> \
    --num_states 3 --num_em_iters 4 --m_epochs 16 \
    --time_range_sec 0 80
\end{verbatim}

\item \textbf{Fit with Lasso-Poisson}:
\begin{verbatim}
./fit_lassoPoisson3.py --basePath $basePath \
    --dataName <dataName> --num_epochs 300
\end{verbatim}

\end{enumerate}

% ===============================================================
\section{Key Design Choices}
\label{sec:design}
% ===============================================================

\paragraph{Shared $A_{\text{true}}$ across all states.}
The synaptic weight matrix is identical in every state; only the
bias vector $B_m$ differs between states.
This models a circuit whose connectivity is fixed but whose
input drive (e.g.\ neuromodulatory tone, sensory stimulus level)
shifts the population-wide activity baseline between states.
The inference task is therefore to recover both $A_{\text{true}}$
and $\{B_m\}$ from a single non-stationary spike record.

\paragraph{Row convention for Dale's principle.}
Row $i$ of $A_{\text{true}}$ contains the outgoing weights of
neuron $i$.
Excitatory neurons (rows $0 \dots N_E - 1$) have all non-zero
weights $> 0$; inhibitory neurons (rows $N_E \dots N-1$) have
all non-zero weights $< 0$.
This guarantees that neuron $i$ can only excite \emph{or} only
inhibit all of its postsynaptic partners, never both.

\paragraph{Spectral radius and stability.}
$\rho(A_{\text{true}}) = R < 1$ is a sufficient condition for
the stationary Poisson VAR(1) to admit a well-defined stationary
distribution.
In the non-stationary regime, stability is maintained instantaneously
because $A_{\text{true}}$ is unchanged; only the bias shifts.

\paragraph{Smooth simplex interpolation.}
The clipped gradient update (Eq.~\eqref{eq:smooth_update}) prevents
abrupt discontinuities in $B_{\mathrm{eff},t}$ at state-boundary
bins.
This mimics a biological circuit that cannot switch network state
instantaneously.
The mixing speed $\nu$ is independent of the mean dwell time and
can be set to match a desired biophysical transition timescale.

\paragraph{Oracle score as a decoder upper bound.}
The oracle uses exact ground-truth parameters and per-bin
maximum-likelihood, achieving the best possible state assignment
from the spike observations.
Any EM or variational algorithm must approach, but cannot exceed,
this score.
The oracle score therefore quantifies how much information about
the latent state is present in the spikes at all, distinguishing
fundamental ambiguity from algorithm suboptimality.

\end{document}
